%---------------------------------------------------------------
\chapter*{Úvod}\addcontentsline{toc}{chapter}{Úvod}\markboth{Úvod}{Úvod}
%---------------------------------------------------------------
\setcounter{page}{1}

% Done - Introduction why this project is needed
V posledních letech dochází k výrazným změnám v oblasti vzdělávání, a to zejména díky rozvoji nových technologií a metod učení. Stále větší popularitu získává mikrolearning, tedy učení prostřednictvím krátkých a efektivních lekcí. Ten však zatím není ve většině oborů dostatečně rozšířen. S nástupem umělé inteligence, zejména velkých jazykových modelů, se také zásadně mění možnosti, jak efektivně a personalizovaně přistupovat k učení. Tradiční systémy, jako jsou například vzdělávací kurzy nebo aplikace využívající spaced repetition algoritmy, často nedokáží držet krok s tímto rychlým vývojem. V moderních aplikacích je navíc kladen čím dál tím větší důraz na uživatelskou zkušenost, která stále u mnoha vzdělávacích systémů není na dostatečně vysoké úrovni.

% Done - Main goal of the project
Na základě těchto skutečností jsme se rozhodli s kolegou Bc. Jakubem Vondráčkem vytvořit moderní vzdělávací systém ve formě webové aplikace, který bude reflektovat aktuální změny v oblasti vzdělávání. Hlavním cílem této aplikace je propojit efektivní metody učení, s kvalitní uživatelskou zkušeností a možnostmi, které poskytuje umělá inteligence. V rámci této práce bude aplikace zaměřena především na výuku programování a softwarového inženýrství. Bude ovšem navržena tak, aby byla snadno rozšiřitelná a využitelná v různých oblastech vzdělávání. Důraz bude kladen na zábavnost, efektivitu a personalizaci učení, které jsou klíčové pro dlouhodobou motivaci uživatelů.

% Write - Rozšíření existujícího systému
Systém bude rozšířením již existující aplikace na výuku jazyků, kterou jsme společně s kolegou Bc. Janem Hlaváčem vytvořili v rámci našich bakalářských prací. Díky tomu bude možné projekt stavět na již vybudovaných základech a zaměřit se tak na implementaci pokročilejších funkcionalit. % Přidat citaci na moji a Honzovu práci? Dozjistit, jestli se může v úvodu citovat - nejspíš ano

% Done - Division into two master theses
Projekt byl rozdělen do dvou diplomových prací. Tato práce bude zaměřena na podrobnější analýzu požadavků a následný vývoj frontendové části aplikace. Diplomová práce kolegy Bc. Jakuba Vondráčka se pak bude převážně zabývat vývojem backendové části této aplikace.

% Done - Subgoals
Z hlavního cíle vytvoření vzdělávací aplikace přirozeně vyplývají dílčí cíle, které reflektují jednotlivé fáze životního cyklu softwarového projektu. Prvním cílem je provést analýzu existující aplikace na výuku jazyků a sepsat požadavky, které budou sloužit jako podklad % Zde nevím, jestli podklad je vhodné slovo, ale asi ano
pro vytvoření nové aplikace. V rámci analýzy bude také provedena analýza konkurenčních řešení v oboru výuky softwarového inženýrství. Druhým cílem je sepsat specifikaci nových funkcionalit a navrhnout jejich uživatelské rozhraní. Součástí tohoto cíle je také provedení aktualizace vzhledu existujícího uživatelského rozhraní. Třetím cílem je provedení samotné implementace funkcionalit a popsání použitých postupů a nově zavedených technologií. Čtvrtým cílem je zvolit a provést vhodné formy testování nových funkcionalit. Posledním cílem je pak zhodnocení celého řešení a navrhnutí možných úprav do budoucna.

% Done - Disclaimer
Práce bude zaměřena výhradně na vývoj aplikace a nebude zahrnovat vytváření samotného obsahu, jako například tvorbu konkrétních kurzů, lekcí nebo cvičení.

% Done - How the project will benefit others
Výsledná aplikace bude sloužit studentům a samoukům, kteří chtějí rozvíjet své znalosti a systematicky se zdokonalovat napříč různými obory. Výstup této práce bude mít největší přínos pro uživatele, kteří se chtějí vzdělat v oblastech programování a softwarového inženýrství, zároveň díky její rozšiřitelnosti bude možné aplikaci využít i v mnoha dalších oborech.
